Abbe number                                          : see Constringence.
Aberration                                           : blemish of an optical system producing blurred or distorted image.
Achromat                                             : optical system correcting chromatic aberration.
Acoustic frequency                                   : frequency of an acoustic wave.
Acoustic velocity                                    : propagation speed of an acoustic wave.
Acousto optic effect                                 : diffraction of light in a material in which propagates an acoustic wave generating periodic variations in refractive index.
Acousto optic modulator                              : optical modulator using acousto optic effect.
Afocal system                                        : optical system without focal planes, so that the image of a point at infinity is also at infinity.
Air corner interferometer                            : interferometer made of two non-parallel optical surfaces separated by air.
Airy pattern                                         : pattern representing the intensity distribution of a plane monochromatic light wave diffracted by a circular aperture.
Albedo                                               : ability of a surface to reflect light defined by the ratio of reflected power to incident power.
Anamorphose                                          : non-homotetic modification of the shape of a light wave.
Angle of incidence                                   : angle of a light ray incident on a surface with the normal to the surface.
Angle of refraction                                  : angle of a light ray refracted by a surface with the normal to the surface.
Angular aperture (in the image space)                : apparent angle of the aperture of an optical system in the image space seen from the on-axis image.
Angular aperture (in the object space)               : apparent angle of the aperture of an optical system in the object space seen from an on-axis object.
Angular deflection                                   : angular deviation.
Anisotropic medium                                   : medium in which the speed of a light wave and therefore the refraction index depend on the direction of propagation and on the polarization.
Anti reflection coating                              : layer(s) deposited on a substrate in order to minimize the reflection of light on the said substrate.
Aperture                                             : diaphragm limiting the number of light rays that are able to propagate.
Apex                                                 : top of a surface.
Attenuation per length                               : expression of the ratio between the intensity of a wave before and after propagation in a medium ( i.e. fiber optic, atmosphere, glass, etc.. ) over a unit of length.
Axial chromatism                                     : wavelength dependance of the position of the paraxial image by an optical system caused by the chromatic aberration for an on-axis object.
Back focal length                                    : for an optical system, algebraic distance from the vertex of the last surface to the back focal point.
Back focal plane                                     : plane conjugate by an optical system of an object plane perpendicular to the optical axis and located at infinity.
Back focal point                                     : conjugate by an optical system of the on-axis point located at infinity.
Back nodal point                                     : point conjugate by an optical system of the front nodal point.
Bandwidth                                            : width ( expressed in frequency or wavelength unit ) of the spectrum in which a phenomenon is supposed to happen.
Birefringence                                        : property of an anisotropic medium in which light propagates differently depending on its polarization.
Birefringent plate                                   : plate made of a birefringent material.
Black body                                           : theoretical object absorbing all the electromagnetic energy it receives, which creates thermal agitation; at equilibrium ( i.e. at constant temperature), all the energy absorbed is re-emitted in electromagnetic radiation also called thermal radiation.
Bragg angle ( acousto optic )                        : in Bragg configuration, angle between the incident rays and the equiphase planes of the ultrasonic wave propagating in the acousto optic crystal ; this angle is the opposite of the angle that the rays diffracted by the acousto optic crystal make with the equiphase planes.
Bragg configuration (acousto optic)                  : configuration in which an incident monochromatic light wave is mainly diffracted in the first order with a direction corresponding to the specular reflection by equiphase planes of the ultrasonic wave.
Brewster angle                                       : particular angle of incidence of a light wave on a refracting surface annihilating? the reflection for the TM polarization.
Bright fringe                                        : said of a constructive interference fringe.
Catadioptric system                                  : optical system made of reflective and refractive optical components.
Cavity mode                                          : resonant wave of a cavity.
Center thickness                                     : for an optical component or system centered on an optical axis, distance between the vertexes of the first and the last optical surface.
Chromatic aberration                                 : see chromatism.
Chromatic dispersion ( fiber optic )                 : duration broadening of a light pulse propagating in a fiber optic and caused by the dependance on the wavelength of the group velocity.
Chromatic dispersion factor ( fiber optic )          : expresses the chromatic dispersion of a light pulse with a 1 nm linewidth and propagating in a 1 km long fiber optic.
Chromatism                                           : aberration occurring in optical systems with refractive surfaces and caused by the dependence of the refractive index on the wavelength.
Cladding ( fiber optic )                             : part of a fiber optic surrounding the core.
Clear aperture                                       : aperture diameter of an optical system.
Coherence length                                     : length characterizing the maximum distance between photons emitted by the same light source so that they can interfere.
Coherence time                                       : time interval characterizing the maximum delay between the emission of two photons by a single light source so that they can interfere.
Collimated beam                                      : beam made of parallel light rays.
Confocal cavity                                      : cavity formed by two coaxial mirrors whose focal points are coinciding.
Conjugation                                          : correspondence between two entities, one being the image of the other ( the object ) through an optical system.
Constringence                                        : parameter characterizing the chromatic dispersion of a material - the higher the constringence, the lower the dispersion.
Core                                                 : central part of a fiber optic in which lightwaves are guided.
Cosine type interferometer                           : Two-wave interferometer in which the phase shift between the interfering waves is a cosine function of the difference in optical paths.
Coupling efficiency ( fiber optics )                 : ratio between the optical power emitted by a source and the optical power effectively injected in a fiber optics.
Coupling ratio ( fiber optics )                      : see « coupling efficiency ».
Critical interaction length ( acousto optic )        : for a monochromatic light wave propagating perpendicularly to the acoustic wave, distance covered for running through half of an acoustic wavelength in the transverse plane ( plane perpendicular to the light wave propagation ) because of the diffraction - defines an interaction length far below which an acousto optic material may diffract in Raman Nath regime.
Cutoff frequency                                     : frequency above which a device is not considered to perform adequately.
Cutoff wavelength                                    : wavelength above which a fiber optic is single mode and below which it is multimode.
Cylindrical lens                                     : lens made of two cylindrical surfaces with parallel axis.
Dark current                                         : electric current flowing through a photosensitive device in complete darkness (in the absence of incident light).
Dark fringe                                          : destructive interference fringe.
Defocus                                              : considering a focus plane conjugate of an object plane by an optical system, defines a shift of the focus plane from its best position (for which the image is the sharpest).
Depth of field                                       : considering a focus plane conjugate of an object plane by an optical system, defines the maximum shift of the object plane enabling a sharp enough image in the fixed focus plane.
Depth of focus ( laser beam )                        : for a laser beam, defines the maximum length from the waist and along the propagation axis of a portion of beam with a size smaller than a given value - for gaussian beams, this value is often square root of 2 times the waist size.
Depth of focus ( optical system )                    : for an optical system with an optical axis, defines the maximum translation on the optical axis of the image plane giving a sharp enough image with a fixed object plane.
Detectivity ( photodetector )                        : inverse of noise equivalent power ( NEP ).
Diffraction                                          : divergent nature of a wave due to its limited size.
Diffraction angle                                    : angle of a diffracted light ray with the normal to the diffracting surface.
Diffraction efficiency                               : for a diffracted light wave (diffraction by a grating, an acousto optic crystal or whatever diffracting component), ratio of the diffracted wave intensity over the incident wave intensity.
Diffraction limited                                  : said of an optical system with no aberration or aberration effects negligible regarding the diffraction effect.
Diffraction order                                    : for a diffraction grating or an acousto optic crystal and a given monochromatic plane light wave, defines the difference between the optical paths expressed in wavelength unit and traveled by light rays diffracted by two consecutive diffractive elements ; also defines the directions along which the diffracted wave may propagate.
Diffusion                                            : phenomenon by which a radiation is deflected in various directions by interaction with other objects.
Dioptre                                              : spherical surface interfacing two transparent media.
Dispersion ( grating )                               : for a given monochromatic plane wave diffracted by a grating in a certain order, defines the angle tilt of the diffracted wave for a wavelength change of 1 nm.
Dispersion ( refractive material )                   : measure relative to the dependance of the refractive index of a material on the wavelength.
Dispersion of deviation ( prism )                    : difference between the deviations of a prism at two different wavelengths.
Divergence                                           : defines the maximum angle between the light rays of a beam.
Diverging light source                               : light source emitting a light beam with a non null divergence.
Doppler Fizeau effect                                : frequency shift between emitted and received light waves caused by the relative speed between the emitter and the receiver.
Duty cycle                                           : for a pulsed periodic light source, especially a laser, defines the ratio between the pulse duration and the period.
Dynamic contrast ( acousto opic )                    : for an acousto optic modulator, contrast of the modulated diffracted beam (limited by the modulation frequency).
Effective focal length                               : algebraic distance from the primary/secondary principal plane of an optical system to its front/back focal plane.
Efficiency of an acousto optic                       : percentage of incident light intensity transferred into the diffracted wave.
Electronvolt                                         : kinetic energy acquired by an electron from zero speed and submitted to a voltage of one volt.
Elliptical beam                                      : a beam in which « equi-intensity » curves in a plane perpendicular to the propagation axis are elliptical.
Ellipticity factor                                   : for an elliptical beam, ratio between the axis lengths of the ellipse containing most of the intensity in a plane perpendicular to the propagation axis.
Embedded gaussian beam                               : beam with a gaussian intensity distribution used to simulate a real beam from the gaussian fundamental mode ; characterized by the M-square coefficient representing the ratio between its divergence and that of the fundamental gaussian mode with the same waist size.
Emissivity                                           : ratio between the emittance of a light source and that of a black body at the same temperature.
Emittance                                            : optical power per unit area emitted by a light source.
Energy per pulse                                     : for a pulsed light source, optical energy emitted in a single pulse.
Equal inclination interferometer                     : two waves interferometer generally formed by two parallel reflecting surfaces so that the phase shift between interfering monochromatic waves is constant with the angle of incidence.
Equiphase surface (interferometer)                   : surface where the phase shift between interfering waves is constant.
Equiphase surface (wave)                             : for a light wave, surface where the phase is constant.
Equivalent mirror (catadioptric system)              : in Gauss approximation, single mirror behaving in the same way as the catadioptric system.
Etendue                                              : describes the envelope in which the light from a source passing through an aperture is wrapped; for a small source and a small aperture in air or vacuum, product of the apparent area of the source seen from the aperture by the solid angle of the aperture seen from the source ; in an equivalent way, product of the apparent area of the aperture seen from the source by the solid angle of the source seen from the aperture.
Extraordinary refraction index                       : in a birefringent uniaxial material, said of the refractive index seen by that of the two refracted rays from a single incident ray whose polarization is not perpendicular to the optical axis.
F-number                                             : dimensionless number whose inverse value is twice the numerical aperture of an optical system ; for optical systems designed to « work » at infinity, it is commonly assumed as the ratio between the focal length and the aperture diameter.
Fabry Perot cavity                                   : resonating cavity made of two reflecting surfaces.
Fabry Perot interferometer                           : multi-wave interferometer made of a parallel plate which surfaces have a non null reflectivity.
Fast axis                                            : for a birefringent medium and a given direction of light wave propagation, said of the polarization orientation with the lowest refraction index.
Fiber optic                                          : long cylindrical optical component guiding light.
Field                                                : zone in the object plane effectively conjugate by an optical system without major vignetting.
Finesse of a Fabry Perot                             : coefficient characterizing the spectral selectivity of a Fabry Perot and approximately equal to the ratio between the free spectral range and the linewidth of a resonant mode.
Focus ( to )                                         : make light rays converge into a single point or diverge from a single point.
Fraunhofer approximation ( Fraunhofer diffraction )  : approximation used for the calculation of the diffracted wave and relevant when the diffractive element and the diffracted pattern are small enough regarding their distance.
Fraunhofer diffraction                               : diffraction of a monochromatic light wave by a plane diffractive element (aperture, etc...).
Free spectral range ( Fabry Perot )                  : distance between two consecutive longitudinal modes.
Free spectral range ( grating )                      : maximum spectral width making it possible to avoid the overlap of a diffraction order with its neighbors.
Frequency                                            : reverse value of the period.
Frequency bandwidth                                  : see bandwidth.
Fringe ( see Interference fringe )                   : bright or dark band.
Fringe order                                         : for two monochromatic and isochromatic waves interfering, ratio between the optical paths difference and the wavelength - fringe order is k for constructive interference ( bright fringes ) and k+ 0,5 for destructive interference ( dark fringes ), k being an integer.
Front focal length                                   : for an optical system, algebraic distance from the vertex of the first surface to the front focal point.
Front focal plane                                    : plane perpendicular to the optical axis whose image plane by the optical system is located at infinity.
Front focal point                                    : point on the optical axis whose image by an optical system is located at infinity.
Front nodal point                                    : For an optical system with a symmetry of revolution around an optical axis, intersection of the primary principal plane with the optical axis.
Full well capacity ( image sensor )                  : maximum charge that a single pixel can contain before saturating.
FWHM                                                 : full width half maximum.
Gain ratio                                           : for any amplifying component or system, defines the ratio between output power and input power.
Gauss condition                                      : applies to an optical system when the object is close to the optical axis and the ray direction is close to that of the optical axis, so that the paraxial approximation can be used.
Gaussian beam                                        : beam with a gaussian distribution of intensity in the plane perpendicular to the direction of propagation - fundamental transverse mode of a laser.
Grating                                              : optical device made of periodically spaced out diffractive elements.
Grey body                                            : body emitting radiations in constant proportion to the black body at the same temperature but with an emissivity lower than 1.
Height chromatism                                    : for an off-axis object, wavelength dependance of the height of the paraxial image by an optical system caused by the chromatic aberration.
Illuminance (photometry)                             : luminous flux received per unit area (as perceived by the human eye and therefore expressed in photometric unit ); equivalent to irradiance in radiometry.
Image                                                : conjugate of an object by an optical system (see object).
Image space                                          : space conjugate of the object space by an optical system.
Image spot                                           : spot described by the intersection of the observation plane with light rays coming from an object point and transmitted through an optical system.
Integrating sphere                                   : spherical cavity with a white and uniformly diffusing coating and therefore behaving as a lambertian diffuser.
Intensity ( photometry )                             : luminous flux emitted by a light source per unit of solid angle (as perceived by the human eye and therefore expressed in photometric unit); same term used in radiometry.
Intensity ( physical optics )                        : local parameter proportional to the average energy or power of a light wave per unit area.
Intensity ( radiometry )                             : power emitted by a light source per unit of solid angle ; same term used in photometry.
Intensity profile                                    : for a light beam, intensity distribution along an axis perpendicular to the propagation axis.
Interaction height ( acousto optic )                 : in an acousto optic crystal, defines the length of the zone parallel to the direction of propagation of the acoustic wave in which the interaction between acoustic and light waves occurs.
Interaction length ( acousto optic )                 : in an acousto optic crystal, defines the length of the zone perpendicular to the direction of propagation of the acoustic wave in which the interaction between acoustic and light waves occurs.
Interference                                         : superposition of light waves whose resulting intensity depends on the phase shifts between them.
Interference fringe                                  : bright or dark band caused by light waves interference - bright and dark fringe occur for waves respectively in phase and out of phase with one another.
Interferometer                                       : device generating interference.
Interfringe                                          : distance between two consecutive bright or dark fringes of an interference pattern.
Irradiance ( radiometry )                            : optical power received by unit area.
Isotropic light source                               : light source emitting with a constant intensity whichever the direction.
Johnson noise                                        : noise generated in an electronic device by the thermal agitation of the carriers (electrons).
Lambertian diffuser                                  : surface diffusing light uniformely whichever the direction and therefore behaving as a lambertian light source.
Lambertian light source                              : light source with a constant radiance (or luminance) whichever the direction.
Laser                                                : acronym of                                   : « light amplification by stimulated emission of radiation » ; device emitting spatially and temporally coherent light radiation based on the stimulated emission.
Lateral colour                                       : see Height chromatism.
Lens                                                 : optical component with two refractive surfaces axially symmetrical for redirecting the light in a controlled manner.
Linewidth                                            : width of a narrow spectrum such as a longitudinal mode of a laser or a line of luminescent source ....
Littrow angle                                        : for a reflective grating, angle of incidence for which the diffracted wave in the order 1 propagates in the same direction than the incident wave.
Longitudinal mode                                    : resonnant wavelength or frequency of a laser or a Fabry Perot interferometer.
Longitudinal spherical aberration                    : algebraic distance from the paraxial image to the marginal image obtained with rays propagating through an optical system from an on-axis point.
Loss ratio                                           : for any optical component or system introducing losses, defines the ratio between output power and input power.
Luminance ( photometry )                             : for a light source or a diffusing/reflecting surface, emitted luminous power per unit solid angle and unit area expressed in photometric unit.
Luminous efficiency                                  : for a light source, ratio between the luminous flux over the power consumption ( expressed in Lumen/W ).
Luminous flux                                        : characterizes the power of a considered light wave as perceived by the human eye and therefore expressed in photometric unit.
M square factor (M2)                                 : ratio between the divergence of an embedded gaussian beam and that of the fundamental mode of a laser (gaussian mode) with the same waist size and at the same wavelength.
Magnification                                        : transverse magnification ( in the calculation and tutorial pages ).
Marginal image                                       : point into which marginal rays, coming from an object point on the optical axis, are converging after being transformed by an optical system.
Maximum exposure time                                : maximum time during which an imaging sensor is exposed to light before reading.
Merit factor ( acousto optic )                       : value depending on some physical parameters of a material and characterizing its acousto optic performances.
Michelson interferometer                             : device splitting a light wave in two perpendicular paths, each split wave being reflected back and recombined with the other for interfering.
Minimum deviation ( prism )                          : for a prism, minimum angle possible between incident and refracted light rays.
Mirror                                               : surface reflecting light rays symmetrically regarding the normal to the surface (specular reflexion) and therefore without scattering them.
Modal dispersion                                     : in a multimode fiber optic, spread in time of a pulse caused by the different propagation velocities of the different modes.
Mode lock                                            : in a laser, longitudinal modes synchronization enabling short pulses emission.
Modulation transfer function (MTF)                   : Performance metric measuring the degradation of contrast in the image compared to contrast in the object, the image being generated by an imaging device including an optical system and/or a camera..
Multimode fiber optic                                : fiber optic enabling the propagation of several transverse modes.
NA                                                   : acronym for numerical aperture.
NEP                                                  : acronym for noise equivalent power.
Newton ring                                          : pattern formed by bright and dark rings resulting from interference obtained from a monochromatic plane wave illuminating a two-wave interferometer made of two reflecting surfaces, one of which is curved and the other is plane.
Noise equivalent power (photodetector)               : minimum detectable power per square root bandwidth – corresponds to the minimum detected optical power generating a photocurrent equal to the RMS noise (Root Mean Square noise) in a 1 hertz bandwidth.
Normal incidence                                     : incidence of light rays when perpendicular to the considered surface.
Normalized frequency ( fiber optic )                 : see V number.
Normalized intensity                                 : ratio of the intensity to its maximum value.
Numerical aperture ( fiber optics )                  : dimensionless number characterizing the acceptance cone angle of an optical fiber.
Numerical aperture ( optical system )                : dimensionless number characterizing the aperture in the object space or in the image space and defined as the product of the refraction index by the sine of half the aperture angle.
Object                                               : element (point, line, plane, or whatever,…) of the space linked to another element called image in a way that rays coming from the object and transmitted by an optical system focus on the image ; object and image are said to be conjugate by the optical system.
Object nodal point                                   : see front nodal point.
Object space                                         : space considered before transformation of light rays by an optical system.
Optical axis                                         : for certain optical systems, axis around which the said optical system has a symmetry of revolution.
Optical axis ( birefringent material )               : direction of propagation of a wave in a birefringent material for which the polarization remains unchanged.
Optical center                                       : for certain optical systems, point of the optical axis coinciding with its conjugate by the optical system.
Optical length of a cavity                           : optical length corresponding to a single round trip in the cavity.
Optical path ( optical length )                      : product of a length (along the path of a light ray) by the refraction index of the medium.
Optical power density ( light source )               : optical power emitted per unit area.
Optical system                                       : in geometrical optics, assembly of optical components for redirecting the light in a controlled manner.
Optoelectronic component                             : component with an optical function derived from electrical energy or vice versa ( laser diode, acousto optic, photodetector, etc.).
Order ( interference )                               : number of wavelengths corresponding to the difference between the optical paths traveled by two interfering light rays.
Ordinary refraction index                            : in a birefringent uniaxial material, said of the refractive index seen by that of the two refracted rays from a single incident ray whose polarization is perpendicular to the optical axis.
Paraxial image                                       : image of an object by an optical system in Gauss conditions.
Peak power                                           : for a pulsed light source, especially a laser, maximum instantaneous power of a pulse.
Period ( of a light wave )                           : time between two consecutive identical states of the electric ( or magnetic ) field.
Phase                                                : characterizes the instantaneous relative amplitude value of a light wave.
Phase dispersion                                     : derived function of the phase with respect to the wavelength.
Phase shift                                          : characterizes the temporal shift between two light waves.
Photocurrent                                         : in a photodetector, current generated by the detected light.
Photodetector                                        : device sensitive to light and converting the absorbed photons into a signal, in most cases into an electrical signal.
Photometry                                           : field of science studying the energetic characteristics of visible light radiation as perceived by the human eye.
Photon noise                                         : intrinsic noise in light caused by its corpuscular nature.
Photopic vision                                      : vision with daytime lighting.
Pitch ( grating )                                    : distance between adjacent diffractive elements of a grating corresponding to the period of the diffracting pattern.
Pixel                                                : smallest sensitive element of an imaging sensor.
Plane parallel plate                                 : optical component made of a refractive medium in between two plane and parallel surfaces.
Polarization                                         : electric field or magnetic field direction of a light wave.
Power ( of an optical system)                        : inverse value of the effective back focal length.
Power density                                        : emitted or received power per unit area.
Power of resolution (optical system)                 : minimum distance ( or minimum angular distance ) between distinguishable points in an image formed by an optical system.
Power of resolution (spectrometer)                   : ratio between the wavelength and the minimum distinguishable variation of wavelength.
Primary principal plane                              : plane of the object space of an optical system whose conjugation by the said optical system occurs with a magnification of +1 and whose conjugate plane is called the secondary principal plane.
Primary principal point                              : intersection of the primary principal plane with the optical axis.
Principal axial chromatism                           : axial chromatism of an optical system for an object at infinity.
Prism                                                : optical component made of a refractive medium in between two unparallel and polished plane surfaces.
Prism angle                                          : angle between the two optical surfaces of a prism.
Quality factor ( acousto optic )                     : parameter determining the interaction regime of a light wave with an acousto optic component crossed by an acoustic wave ; the interaction is the Raman Nath regime when the quality factor is far below 1 and the Bragg regime when it is far above 1.
Quantum efficiency                                   : percentage of photons absorbed by a photodetector and effectively producing a signal.
Radiance ( radiometry )                              : for a light source or a diffusing/reflecting surface, power emitted per unit of solid angle and per unit of apparent area.
Radiometry                                           : field of science studying the energetic characteristics of radiation, including radiation in the optical spectral range.
Raman Nath configuration                             : for an acousto optic crystal, configuration in which a light wave is diffracted in different directions with different frequencies ( or wavelengths ) by a thin acoustic wave propagating perpendicular to it.
Rayleigh range                                       : for gaussian or embedded gaussian beams, distance in the propagation direction from the waist to a position where the beam radius is increased by a factor of square root of 2.
Readout noise ( image sensor )                       : noise introduced by the on-chip amplifier during the conversion of the charge of a single pixel ( electrons successfully hit by photons ) into a signal ( current or voltage ).
Reflectance                                          : reflection coefficient related to the intensity.
Reflection                                           : phenomenon occurring when light hits a material and related to light rays which are neither absorbed nor transmitted by the said material.
Reflection coefficient                               : ratio between reflected light power ( or intensity or amplitude ) and incident light power ( or intensity or amplitude ).
Reflective grating                                   : diffraction grating made of reflective diffractive elements.
Refracting surface                                   : surface interfacing two transparent media with different refraction indexes.
Refraction                                           : deviation of rays at the interface of two transparent media with different refraction indexes according to Snell-Descarte law.
Refraction index                                     : for a transparent medium, ratio between the light speed in the air and the light speed in the medium ; for a reflective surface in the air or vacuum, the refraction index is « - 1 ».
Relative intensity                                   : normalized intensity as determined by the ratio between the intensity and its maximum value.
Repetition rate                                      : number of time per second that a pulse is emitted by a light source.
Resolvance (grating)                                 : ratio between the considered wavelength and the smallest resolvable wavelength variation ; two wavelengths are considered resolved if the distance between their diffraction peaks is at least equal to half the width of a diffraction peak.
Resonant Cavity                                      : cavity in which only some specific light waves may oscillate, possibly oscillating light waves being the resonant modes of the cavity.
Resonant wavelength of a Fabry Perot                 : see longitudinal mode.
Rigorous stigmatism                                  : Property of certain optical systems for certain characteristic points of the object space such that all the rays coming from one of these points of the object space rigorously pass through a point of the image space (conjugate point) after transmission by the system optical.
Rise time ( acousto optic )                          : time it takes for the acoustic wave to run through the beam width.
Schmidt plate                                        : refractive plate perpendicular to the optical axis compensating the spherical aberration of a spherical mirror when the object is at infinity.
Secondary light source                               : object diffusing and eventually reflecting / transmitting the light coming from a light source generally called primary light source.
Secondary principal plane                            : plane conjugate by an optical system of an object plane ( the primary principal plane ) with a transverse magnification of +1.
Secondary principal point                            : intersection of the secondary principal plane with the optical axis.
Sensitivity                                          : for a photodector, ratio between the photocurrent and the incident optical power.
Shot noise                                           : electronic noise caused by the random motion fluctuations of charge carriers (electrons) in a conductor; in a photodetector, the shot noise is partly intrinsic and partly coming from the photon noise.
Signal to noise ratio                                : (often called by its acronym                                   : SNR) ratio between the signal and the noise; for a photodetector, the signal is the current or voltage generated by the detected light.
Single mode fiber optic                              : fiber optic enabling the propagation of only one transverse mode ( the fundamental mode ).
Slow axis                                            : for a birefringent medium and a given plane light wave, polarization direction with the highest refraction index.
SNR                                                  : see signal to noise ratio.
Solid angle                                          : two-dimensional angle in the three-dimensional space measuring how large an obect appears from a viewing point.
Spatial frequency                                    : characterizes a periodic structure in a direction of the space; gives the number of periods of the structure per unit length.
Spatial mode                                         : specific spatial distribution of the electromagnetic field in a plane perpendicular to the propagation axis of a light wave possibly emitted by a laser or possibly propagating in a waveguide ( like, for instance, a fiber optic ) - also called transverse mode.
Specific detectivity                                 : product of the detectivity by the square root of the area of the photodetector.
Spectral power density                               : power density per unit wavelength.
Spectral sensitivity ( human eye )                   : normalized sensitivity of the human eye to light depending on the wavelength; calculated in the visible spectral range with a maximum value at 555 nm for the photopic vision.
Spectral width                                       : difference between the largest and smallest wavelengths (frequencies) of the spectrum emitted by a light source, transmitted by an optical system or detected by a photodetector, etc.; the largest and smallest wavelengths (frequencies) are often defined at Full Width Half Maximum (FWHM).
Spectrum                                             : contribution of the monochromatic components of a wave on a phenomenon (emission, detection, absorption, transmission, etc.), as a function of the wavelength (or frequency) of its components.
Specular reflexion                                   : property of the surface of a material with low rugosity reflecting light without scattering ; the specular reflexion redirects an incident light ray in a reverse way and in a direction symmetrical to the incident ray regarding the normal to the surface.
Spherical aberration                                 : geometric aberration induced by an optical system on rays from on-axis objects.
Spherical lens                                       : lens made of two coaxial spherical surfaces.
Spherical mirror                                     : mirror with a spherical reflecting surface.
Stability coefficient ( laser cavity )               : parameter depending on radii of curvature of the mirrors and on the distance between mirrors and which value is between 0 and 1 when transverse modes of a cavity may exist.
Stability condition of a resonator                   : condition defining if transverse modes of a cavity may exist or not.
Stigmatic                                            : property of an optical system realizing perfect conjugation between an object and its image.
TE polarization                                      : polarization of a wave whose electric field is perpendicular to a considered plane ; in the case of a wave incident on a surface, the considered plane is the incidence plane ( containing the wave direction and the normal to the surface ).
Thick lens                                           : lens with a non-negligible center thickness.
Thin lens                                            : lens with a negligible center thickness.
Tilt angle                                           : angle of an entity or direction with a reference axis.
TM polarization                                      : wave whose magnetic field is perpendicular to a considered plane ; in the case of a wave incident on a surface, the considered plane is the incidence plane ( containing the wave direction and the normal to the surface ).
Total deviation ( prism )                            : angular deviation of a light ray after refraction by a prism.
Total reflection                                     : reflection of light by a surface separating two transparent media and during which all the light is reflected; phenomenon only possible if the refractive index of the incident medium is larger than that of the emerging medium and if the angle of incidence is larger than a certain value.
Transmission                                         : ratio between transmitted radiation power (or intensity or amplitude) and incident radiation power (or intensity or amplitude).
Transmission coefficient                             : ratio between transmitted radiation power (or intensity or amplitude) and incident radiation power (or intensity or amplitude).
Transmission grating                                 : diffraction grating made of transmitting diffractive elements.
Transmission ratio                                   : ratio between transmitted radiation power (or intensity or amplitude) and incident radiation power (or intensity or amplitude).
Transversal spherical aberration                     : diameter of the spot formed in the paraxial image plane with rays propagating through an optical system and coming from an on-axis object.
Transverse magnification                             : size ratio of an image and an object, the image being conjugated to the object by an optical system; for centered optical systems, the transverse magnification is defined with objects in a plane perpendicular to the optical axis - transverse magnification is then positive if object and image have the same orientation and negative otherwise.
Transverse mode                                      : see spatial mode.
Uniaxial material                                    : said of a birefringent material for which the ordinary and extraordinary refraction indexes are identical for plane waves propagating in a given direction called the optical axis.
V number                                             : parameter depending on the wavelength and on the characteristics of a fiber optic, enabling to determine if the fiber optic is single-mode or multimode.
Vertex                                               : top of a surface.
Vignetting ( optical system )                        : darkening at the periphery of the image surface.
Waist                                                : position along the propagation axis where a laser beam has the smallest size.
Wavefront                                            : equiphase surface of a light wave.
Wavefront radius                                     : radius of curvature of the wavefront of a spherical wave.
Wavelength ( of a light wave )                       : distance between two consecutive identical states of the electric ( or magnetic ) field.
Wavenumber ( spectroscopy )                          : inverse wavelength value; currently used in spectroscopy instead of the wavelength.
Weierstrass points                                   : for an elliptic mirror or a spherical refracting surface, points of the optical axis which are rigorous conjugates of each other; for an elliptical mirror, the Weierstrass points are the foci of the ellipse.
